\section{Document extraction: the lexical ESG-zone filter}
\label{sec:extraction}

The original study vectorises each report in its entirety: a standalone
sustainability report is, by construction, already ESG text, so
\citet{lagasio2024} needs no step to separate ESG content from anything else.
Our corpus is different by design (\S\ref{sec:scope}, \S\ref{sec:data}): annual
reports and URDs in which ESG disclosure is embedded inside a much larger
financial and legal document. Feeding the whole report to the downstream
measures would let financial-statement bulk and document length dominate
every score. Extraction is therefore a required stage with no counterpart in
the original, and it fixes the text, and hence the denominators, on which every
component of \S\ref{sec:index} is computed.

\subsection{Paragraph segmentation}

The extractor is the class \texttt{LexicalDocumentFilter} in
\nolinkurl{src/lexical_document_filter.py}; its full parameter and dependency
listing is in \S\ref{sec:repro}. Each PDF is read with PyMuPDF's block-level
extractor (\texttt{page.get\_text("blocks")}), text blocks only, sorted
top to bottom and left to right. A block is accepted as a paragraph only
after passing five filters, in order: (i) \emph{margin strip} — dropped if entirely in the top
or bottom margin band ($y_1 < 0.08h$ or $y_0 > 0.92h$, $h$ = page height);
(ii) \emph{running header/footer suppression} — every margin-band block of
at most 120 characters is normalised (digits $\to$ \texttt{\#}) and its
occurrences counted over the document, and any normalised string whose
occurrences exceed 30\% of the page count is blacklisted and dropped wherever
it occurs; (iii)
\emph{financial-table detection} — a block is dropped as tabular only if it
contains \emph{no} ESG keyword (the regex defined below) \emph{and} either
its digit-to-character ratio exceeds 0.35 with at most 6 lines, or every
line is under 60 characters, there are at least two lines, and the first has
no terminal punctuation (period, exclamation mark, question mark, or colon);
the keyword veto keeps a table row that happens to carry ESG content rather
than discarding it as noise; (iv) \emph{section headers} — an all-caps block
(ignoring non-letters), at least two letters, at most 10 tokens, fewer than 60\% of
them with one letter or none, and containing no filename-like string (a
word, a period, then 2--6 word characters, e.g.\ \texttt{report.pdf}), is
emitted as a \texttt{\S~TITLE~\S} marker that tags
subsequent paragraphs with their section, rather than entering the
paragraph stream; (v) \emph{navigation} — a block that survives (i)--(iv) is
cleaned (below) and then dropped if it reads as a table of contents,
a GRI/ESRS correspondence table or a stray navigation line
(\S\ref{sec:extraction-nav}).
Cleaning strips URLs, rejoins hyphenated line breaks, drops non-ASCII
characters, and collapses whitespace. Non-ASCII characters are deleted, not
transliterated, so a term written with a typographic hyphen fuses
(\emph{netzero}, which the regex below does not match) and a word set with a
ligature loses letters; the matches lost this way are not counted.

\subsection{Keyword matching}
\label{sec:extraction-regex}

Paragraphs are scored against a single compiled regular expression,
\texttt{\_ESG\_KW\_RE}, built at import time from the four exported
vocabulary files of \S\ref{sec:lexicon}. Under the adopted specification it
combines 402 natural-form terms (367 from \texttt{esg\_terms.txt} and 35 from
\texttt{esg\_terms\_sectorial.txt}) with 32 patterns (31 from
\texttt{esg\_patterns.txt}, one from \texttt{esg\_patterns\_sectorial.txt});
without the sectoral files it would hold 367 and 31. Each term is converted
mechanically: its words are joined by \texttt{[\textbackslash s\textbackslash
-]+}, so spaces and hyphens are interchangeable (\emph{net zero} matches
\emph{net-zero}); an optional \texttt{s} after the last word covers
\emph{-s} plurals only, so \emph{data breaches} is missed and
\emph{fatalities} is caught only because a pattern (\texttt{fatalit\textbackslash
w+}) supplies it. Terms are ordered by number of words, then by length, both
descending, because alternation is leftmost-first rather than longest-match:
without the ordering \emph{carbon} would shadow \emph{carbon footprint}.
The 32 patterns are read from file as written, not generated, and follow the
terms in the alternation: morphological families
(\texttt{sustainab\textbackslash w+}, \texttt{recycl\textbackslash w+}),
digit-bearing forms (\texttt{co2e?}, \texttt{scope\textbackslash s*[123]},
ISO codes), collocations that fix the ESG sense of a polysemous head
(\emph{social}, \emph{water}, \emph{energy}), and phrases that do not
survive lemmatisation (\emph{paris agreement}, \emph{just transition}). The
whole alternation is wrapped in word boundaries and matched
case-insensitively. A match consumes its span, so \emph{carbon footprint}
counts once, not also as \emph{carbon}.

\subsection{Navigation filter}
\label{sec:extraction-nav}

Tables of contents and GRI/ESRS correspondence tables are lists of section
names, dense in ESG keywords, so the density rule below would select them as
ESG text; \S\ref{sec:lexicon-nav} gives the reason for filtering the text
rather than removing the section-naming terms from the vocabulary. The filter
(\texttt{\_is\_navigation}) acts on the cleaned block and counts
\emph{navigation markers} among its whitespace tokens: a number of one to
three digits preceded by a word and followed by nothing, a capitalised token
or a section number (a page reference between titles); a multi-level section
number (\texttt{2.1}, \texttt{1.3.1}) followed by a capitalised token; a
GRI-style (\texttt{205-3}) or ESRS-style (\texttt{S1-14}) disclosure code; or
\emph{p.}, \emph{pp.}, \emph{page} or \emph{paragraph} followed by a digit. A
block is dropped if it contains two or more dotted-leader page references
(five or more dots, then a page number), or if it has at least four markers
making up at least 8\% of its tokens. A short block, of at most 15 tokens,
is kept if it ends in sentence punctuation (\texttt{.!?:;}) and contains no
leader (four or more underscores or dots); otherwise it is dropped if it
opens with a multi-level section number or a Roman numeral and period, ends
in a one- to three-digit number after a word, or contains a leader. The last
rule also drops short unpunctuated labels such as \emph{Scope 3} standing
alone. Unnumbered headings and glossaries are not targeted.

The thresholds above were calibrated on a hand-labelled sample of keyword
occurrences that is not part of the repository. The constants are therefore
documented but their calibration cannot be reproduced, and the pipeline does
not log how many blocks the filter drops, so its effect on the extraction is
unmeasured.

\subsection{Density scoring and zone construction}

Each paragraph $p$ is scored by keyword density,
\begin{equation}
\mathrm{density}(p) = 100 \times \frac{\mathrm{kw\_count}(p)}{\mathrm{words}(p)},
\label{eq:kw-density}
\end{equation}
where $\mathrm{kw\_count}(p)$ is the number of matches of
\texttt{\_ESG\_KW\_RE} on the cleaned, unlemmatised paragraph text and $\mathrm{words}(p)$ its
whitespace-token count.
A paragraph is \emph{hot} if $\mathrm{density}(p) \geq 1.0$. Every hot
paragraph is expanded to a window of one paragraph either side; overlapping
or adjacent windows ($\mathrm{lo} \leq \mathrm{prev\_hi}+1$) are merged into
one contiguous zone; merged zones under 150 characters are dropped. Zone
statistics (keyword count, word count, density) are recomputed on each
surviving zone's concatenated text, not summed from its paragraphs; because
the term regex allows any whitespace between words, a phrase can match across
a paragraph break there.
Table~\ref{tab:extraction-params} lists the free parameters governing hot-paragraph
detection and zone construction, none tuned by grid search; the thresholds internal to
financial-table, section-header and navigation detection above are additional
constants, given in prose rather than tabulated. Per PDF, zones are sorted by start offset and joined
into one \texttt{relevant\_text} field, which is all that every later stage
(\S\ref{sec:lexicon} onward) reads.

\begin{table}[tbp]
\centering
\caption{Extraction parameters.}
\label{tab:extraction-params}
\begin{tabular}{lr}
\toprule
Parameter & Value \\
\midrule
Keyword-density threshold (hot paragraph) & 1.0 / 100 words \\
Context window & $\pm 1$ paragraph \\
Minimum zone length & 150 characters \\
Header/footer repeat threshold & $>30\%$ of pages \\
Margin bands & top 8\% / bottom 8\% of page height \\
\bottomrule
\end{tabular}
\end{table}

\subsection{Yield}
\label{sec:extraction-yield}

Extraction runs before deduplication, on all 344 PDFs, and every PDF yielded
at least one zone. The extra file is one 2019 report filed twice (a
\texttt{\_repaired} copy), removed in \S\ref{sec:data}; its extraction is
identical to the original's. Table~\ref{tab:extraction-yield} gives the
totals on both bases; every per-document statistic, the yield and the
by-year figures are on the 343 unique reports, as is everything downstream of
preprocessing. Zone and word counts come from \texttt{scripts/analysis\_v3.py};
keyword matches and extraction-stage density from
\texttt{scripts/extraction\_stats\_v3.py}, which applies the extractor's own
compiled regex to each report's \texttt{relevant\_text}. Summing the
per-zone counts instead gives 50 fewer matches (1,071,489 on 343), phrases
that span two concatenated zones.

Yield is extracted words over the words of the full PDF text layer
(\texttt{page.get\_text()} over every page), both counted as whitespace
tokens (\texttt{str.split}), on all 343 reports. The tokeniser is the same;
the text is not. The numerator is cleaned text, and cleaning deletes
standalone non-ASCII tokens (bullets, dashes, currency signs), fuses words
across a non-ASCII hyphen and rejoins line-break hyphenation, all of which
the uncleaned denominator counts. Yield is therefore biased slightly
downward, by an amount the committed scripts do not measure. It is the share
of a report's text layer that survives into its zones, not a share of its
ESG content. The mean report has 389 pages.

\begin{table}[tbp]
\centering
\caption{Extraction volume. Totals on all 344
extracted PDFs and on the 343 unique reports; per-document figures on 343
only. Keyword density is Eq.~\eqref{eq:kw-density} with the full regex on
raw extracted text; it is not \susdens.}
\label{tab:extraction-yield}
\begin{tabular}{lrr}
\toprule
Statistic & 344 PDFs & 343 reports \\
\midrule
Total zones & 66,377 & 66,171 \\
Total extracted words & 28,113,340 & 28,101,967 \\
Keyword matches & 1,072,127 & 1,071,539 \\
Keyword density, aggregate (per 100 words) & 3.814 & 3.813 \\
\midrule
Zones per document (mean, sample sd) & & 192.9, 146.0 \\
Zones per document (min, median, max) & & 9, 162, 1,085 \\
Extracted words per document (mean, median) & & 81,930, 74,802 \\
Words per zone (aggregate) & & 424.7 \\
Words per zone (mean of per-document values) & & 487.2 \\
Keyword density (mean of per-document values) & & 3.673 \\
\midrule
Yield vs.\ full PDF text (mean, median) & & 39.6\%, 38.3\% \\
Yield vs.\ full PDF text (min, max) & & 7.1\%, 75.6\% \\
Yield vs.\ full PDF text (aggregate) & & 39.6\% \\
\bottomrule
\end{tabular}
\end{table}

The aggregate yield is 28,101,967 extracted words out of 70,966,812. By year
(Table~\ref{tab:extraction-by-year}), every volume measure rises and extracted
words grow faster than the full text. Extraction-stage keyword density rises
in every year on both bases, aggregated and as a mean of per-report values
(3.21 in 2018 to 4.52 in 2024).

What these rises do not show is that reports disclose more ESG content.
Yield and extraction density are both functions of the vocabulary: a
paragraph is kept because it is dense in these terms, and density counts the
same terms. The vocabulary was built on reports from across 2018--2024
(\S\ref{sec:lexicon}) and contains the names of reporting regimes
that postdate 2018 (\S\ref{sec:temporal-mandate}), so a shift in reporting
language towards this vocabulary raises both measures whether or not the
underlying disclosure changed, and so does a change in how reports are
organised that makes ESG passages longer and more contiguous. Nor do they
show that extraction behaves alike in every year: nothing in the paper
measures, year by year, how often the vocabulary's matches or the navigation
filter's decisions are wrong. \S\ref{sec:temporal} reads the trends with these
qualifications; here they are recorded as properties of the extraction.

\begin{table}[tbp]
\centering
\small
\caption{Extraction by year, 49 reports per year (343 in all). Means per
report except the last column; yield is extracted over full-text words;
keyword density is the aggregate per 100 extracted words, as in
Table~\ref{tab:extraction-yield}.}
\label{tab:extraction-by-year}
\setlength{\tabcolsep}{4pt}
\begin{tabular}{lrrrrrrr}
\toprule
Year & Zones & Extracted words & Words/zone & Full-text words & \multicolumn{2}{c}{Yield} & KW density \\
\cmidrule(lr){6-7}
 & & & & & mean & median & \\
\midrule
2018 & 157.7 & 56,484 & 405.3 & 176,519 & 31.6\% & 28.7\% & 3.24 \\
2019 & 157.6 & 60,594 & 444.5 & 179,803 & 34.6\% & 32.6\% & 3.30 \\
2020 & 181.2 & 70,143 & 473.7 & 197,275 & 36.5\% & 35.5\% & 3.49 \\
2021 & 189.8 & 80,845 & 483.3 & 206,427 & 40.1\% & 37.9\% & 3.63 \\
2022 & 195.8 & 89,382 & 521.2 & 216,282 & 42.1\% & 39.3\% & 3.75 \\
2023 & 226.1 & 96,162 & 521.5 & 224,549 & 43.6\% & 41.2\% & 4.00 \\
2024 & 242.1 & 119,901 & 561.0 & 247,447 & 48.4\% & 48.4\% & 4.56 \\
\midrule
All & 192.9 & 81,930 & 487.2 & 206,900 & 39.6\% & 38.3\% & 3.81 \\
\bottomrule
\end{tabular}
\end{table}

\subsection{Design rationale}

The filter is a pure function of the PDF and the lexicon files: no network
call, no external API, no random seed on this path. This matters directly
for the determinism claimed in \S\ref{sec:repro}. A semantic filter, selecting
text by its similarity to anchor queries in the embedding space of a hosted
model, would not satisfy that claim: a paid, versioned, externally hosted
model can change or be deprecated independently of this codebase, and the
same PDF is then not guaranteed to yield the same selection on a later run.
A lexical filter also ties extraction to the same vocabulary that \SUS\
counts, which is the design choice argued in \S\ref{sec:lexicon-one}.

\subsection{Recall is bounded by the lexicon}

Extraction and \SUS\ draw on the same vocabulary (\S\ref{sec:lexicon-one}),
but not on the same entries. Every entry of the 399-entry lemmatised \SUS\
vocabulary is the lemma of a term in the extraction regex, and every one of
the 402 terms reaches \SUS: the eleven that preprocessing would empty or
collapse to a generic word are protected as single tokens
(\S\ref{sec:lexicon-represent}), and three pairs of terms share a lemma
(Table~\ref{tab:lexicon-lemma}). The converse fails for the 32 patterns,
which act only at extraction: they can make a paragraph hot and pull a zone
in, but no pattern is itself a \SUS\ entry.
Two bounds follow. A term outside the combined regex triggers no hot
paragraph and enters no zone, so it is invisible everywhere. A pattern can
bring text into a zone, where it feeds \SEN, \QUANT, \HEDGE\ and the length
of the zone, without being counted by \SUS. The two stages also match
slightly different strings for the same term: the extraction regex requires a
space or hyphen between words, the protection of \S\ref{sec:lexicon-represent}
also accepts the fused form (\emph{speakup}), and a lemmatised n-gram also
matches across removed stopwords (\emph{board director} in \emph{board and
directors}). For the heaviest multi-word entry this last difference is 0.3\%
(\S\ref{sec:lexicon-represent}). The extraction side of the coupling is the
larger one: a term in the regex brings further occurrences of itself into the
zones it helps select (\S\ref{sec:lexicon-represent}).
\S\ref{sec:limitations} develops the consequences; they are noted here as
properties of this step.
